\documentclass[12pt,letterpaper]{article}
\usepackage[utf8]{inputenc}

% Ajustamos los márgenes para corregir las advertencias de fancyhdr
\usepackage[letterpaper, margin=1in, bottom=2.5in, top=1.2in]{geometry}

% Cargamos primero el paquete fancyhdr
\usepackage{fancyhdr}
% Corregimos los parámetros headheight y footskip
\setlength{\headheight}{15pt} % Al menos 14.49998pt según la advertencia
\setlength{\footskip}{72pt}   % Al menos 71.60547pt según la advertencia

\usepackage{xcolor}
\usepackage{amsmath}
\usepackage{anyfontsize}
\usepackage{graphicx}
\usepackage{tikz} % Necesario para los elementos gráficos
\usepackage{listings} % Para resaltado de sintaxis
\usepackage{sourcesanspro}
\renewcommand{\familydefault}{\sfdefault}
\usepackage{fontawesome5}
\usepackage{titlesec}
\usepackage{setspace}
\usepackage{hyperref}
\usepackage{mdframed} % Para marcos más estables
\usepackage{qrcode} % Paquete para generar códigos QR

% Cargar tikz y la biblioteca de cajitas para mdframed
\usetikzlibrary{calc,shapes,positioning}

% IMPORTANT: Force the text color to be white for the entire document
\AtBeginDocument{\color{primaryColor}}

% Colors updated with better contrast
\definecolor{myblue}{cmyk}{1.0,0.8,0.05,0.20}
\definecolor{bgColor}{RGB}{15, 22, 36}  % Very dark blue, almost black background
\definecolor{primaryColor}{RGB}{255, 255, 255}  % White text - full white for better contrast
\definecolor{accentColor}{RGB}{255, 212, 59}  % Python Yellow
\definecolor{pythonBlue}{RGB}{120, 180, 255}  % Azul Python más claro para mejor contraste
\definecolor{secondaryColor}{RGB}{220, 230, 240}  % Lighter blue-gray for better contrast
\definecolor{terminalBg}{RGB}{22, 22, 22}  % Terminal background remains dark
\definecolor{terminalFrame}{RGB}{40, 40, 40}  % Terminal frame slightly darker
\definecolor{lineNumberColor}{RGB}{100, 100, 100}  % Color más sutil para los números de línea
\definecolor{dividerColor}{RGB}{60, 70, 90}  % Color más sutil para la línea divisoria

% Colores mejorados para el código con mayor contraste
\definecolor{codeTextColor}{RGB}{255, 255, 255}  % Blanco puro para el texto básico
\definecolor{codeKeywordColor}{RGB}{135, 206, 250}  % Azul cielo para palabras clave
\definecolor{codeCommentColor}{RGB}{144, 238, 144}  % Verde claro para comentarios
\definecolor{codeStringColor}{RGB}{255, 215, 135}  % Amarillo para strings
\definecolor{codeMethodColor}{RGB}{255, 160, 122}  % Naranja claro para métodos
\definecolor{codeFunctionColor}{RGB}{255, 215, 0}  % Dorado para funciones
\definecolor{codeNumberColor}{RGB}{255, 255, 150}  % Amarillo claro para números

% Set page background color
\pagecolor{bgColor}

% Hyperref configuration
\hypersetup{
    colorlinks=true,
    linkcolor=accentColor,
    filecolor=accentColor,
    urlcolor=accentColor,
}

% Configuración básica para listings que evita problemas con saltos de página
\lstset{
  language=Python,
  basicstyle=\normalsize\ttfamily\bfseries\color{codeTextColor}, % Tamaño reducido, en blanco y negrita
  backgroundcolor=\color{terminalBg},
  commentstyle=\color{codeCommentColor},
  keywordstyle=\color{codeKeywordColor},
  stringstyle=\color{codeStringColor},
  numberstyle=\color{lineNumberColor}, % Números de línea más sutiles
  breaklines=true,
  breakatwhitespace=true,
  tabsize=4,
  showstringspaces=false,
  frame=none,
  xleftmargin=15pt,
  xrightmargin=0pt,
  aboveskip=10pt,
  belowskip=10pt,
  numbers=left,
  numbersep=8pt,
  extendedchars=true,
  keepspaces=true,
  columns=flexible,
  lineskip=6pt, % Espacio entre líneas ligeramente reducido
  % Keywords de Python (palabras reservadas)
  morekeywords={self, yield, lambda, with, as, from, True, False, None, import, in, for, 
                if, elif, else, while, return, def, class, try, except, finally, raise, 
                break, continue, global, nonlocal, pass, assert, del, yield from},
  % Funciones incorporadas resaltadas de manera distinta
  emph={[2]print,range,sum,int,str,float,list,dict,set,tuple,next,len,type,map,filter,
         reduce,zip,enumerate,sorted,reversed,min,max,open,any,all},
  emphstyle={[2]\color{codeFunctionColor}}
}

% TERMINAL ESTILO MAC MEJORADA CON BOTONES A LA IZQUIERDA Y MENOS ESPACIO SUPERIOR
\newenvironment{macterminal}{%
    \begin{mdframed}[
        linecolor=terminalFrame,
        backgroundcolor=terminalBg,
        roundcorner=5pt,
        skipabove=10pt,
        skipbelow=10pt,
        linewidth=1pt,
        innertopmargin=10pt, % Reducido
        frametitle={%
            \tikz[baseline=(current bounding box.east), outer sep=0pt]{
                \fill[red!80!black] (0,0) circle (5pt);
                \fill[yellow!80!black] (0.7,0) circle (5pt);
                \fill[green!70!black] (1.4,0) circle (5pt);
            }
        },
        frametitlealignment=\raggedright, % Alineado a la izquierda
        frametitleaboveskip=8pt, % Espacio entre el título y el borde superior
        frametitlebelowskip=0pt, % Reducir el espacio entre el título y el contenido
    ]
    % No necesitamos más espacio vertical aquí
}{%
    \end{mdframed}%
}

% Utilities
\newcommand{\verspace}{\vspace{10pt}}

% Section styling - Esquema más elegante y mejorado contraste
\titleformat{\section}
  {\LARGE\bfseries\color{primaryColor}} % Secciones principales en blanco para mejor contraste y más grandes
  {\thesection. }
  {0pt}
  {}
  []

\titleformat{\subsection}
  {\Large\bfseries\color{accentColor}} % Subsecciones en amarillo Python y más grandes
  {\thesubsection. }
  {0pt}
  {}
  []

\titleformat{\subsubsection}
  {\large\bfseries\color{pythonBlue}} % Subsubsecciones en azul Python más claro y más grandes
  {\thesubsubsection. }
  {0pt}
  {}
  []

% Spacing for sections
\titlespacing*{\section}{0pt}{20pt}{10pt}
\titlespacing*{\subsection}{0pt}{15pt}{7pt}
\titlespacing*{\subsubsection}{0pt}{10pt}{5pt}

% Configure header and footer
\pagestyle{fancy}
\fancyhf{}

% Crear un encabezado elegante para todas las páginas (excepto la portada)
\fancyhead[L]{
    \begin{tikzpicture}[remember picture, overlay]
        % Rectángulo de fondo para la sección Python
        \fill[pythonBlue, rounded corners=3pt] (0,0) rectangle (2.2cm,0.7cm);
        % Texto Python
        \node[text=primaryColor, font=\bfseries] at (1.1cm,0.35cm) {PYTHON};
        
        % Pequeño separador
        \fill[secondaryColor] (2.35cm,0.1cm) rectangle (2.40cm,0.6cm);
        
        % Texto Context Managers con el color de acento
        \node[text=accentColor, font=\bfseries, anchor=west] at (2.35cm,0.35cm) {CONTEXT MANAGERS};
    \end{tikzpicture}
}

% Línea de encabezado
\renewcommand{\headrulewidth}{0pt}

% Línea divisoria más sutil
\renewcommand{\footrule}{
  \vspace{0.5cm}
  \noindent\makebox[\linewidth]{\color{dividerColor}\rule{\linewidth}{0.2pt}}
  \vspace{0.5cm}
}
\renewcommand{\footrulewidth}{0.2pt}
\renewcommand{\footruleskip}{1cm}

% Configure footer with profile info
\fancyfoot[C]{
    \vspace*{0.2cm}
    \noindent
    \begin{minipage}{\textwidth}
        \begin{flushleft}
            % Profile image placeholder with TikZ
            \raisebox{0.7cm}{
            \begin{tikzpicture}[baseline]
                \path[fill=bgColor] (0,0) circle (0.8cm);
                \clip (0,0) circle (0.8cm);
                \node at (0,0) {
                    % Placeholder for profile image
                    \includegraphics[width=1.6cm,height=1.6cm]{profile-image.jpeg}
                };
            \end{tikzpicture}
            }
            \begin{minipage}[b]{0.8\textwidth}
                % Name
                {\large\bfseries\color{primaryColor}Alejandro Sánchez Yalí}
                
                % Professional description
                \par\vspace{1pt}
                {\small\color{secondaryColor}Software Developer | AI \& Blockchain Enthusiast}
                
                % Contact
                \par\vspace{1pt}
                {\small\color{accentColor}\faGlobe\hspace{5pt}\color{secondaryColor}www.asanchezyali.com}
            \end{minipage}
        \end{flushleft}
        \vspace{8pt} % Space at bottom
    \end{minipage}
}

% Estilo para la portada (sin encabezado ni pie de página)
\fancypagestyle{plain}{
    \fancyhf{}
    \renewcommand{\headrulewidth}{0pt}
    \renewcommand{\footrulewidth}{0pt}
}

% Bullet styling for lists
\renewcommand{\labelitemi}{\textcolor{accentColor}{$\bullet$}} % Primer nivel de viñetas en amarillo (mejor contraste)
\renewcommand{\labelitemii}{\textcolor{pythonBlue}{$\circ$}} % Segundo nivel de viñetas en azul

% Aumentar globalmente el tamaño de texto un 20%
\usepackage{relsize}
\AtBeginDocument{\relsize{1}} % Incrementa la fuente en aproximadamente un 20%

% Comando simplificado para la etiqueta Python
\newcommand{\languagetag}[1]{
    \begin{tikzpicture}[baseline]
        \node[fill=pythonBlue, text=primaryColor, rounded corners=5pt, inner sep=7pt] {
            {\normalsize\textbf{#1}}
        };
    \end{tikzpicture}
}

% Alternativa usando una versión modificada de la URL
\newcommand{\elegantqr}[2]{
    \qrcode[height=2.5cm]{#1}
    \\[0.1cm]
    {\hspace{0.2cm}\color{primaryColor}\small #2\par}
}

% Definir el título personalizado con alineación exacta para todos los elementos
\newcommand{\titlepagecontents}{%
    \vspace*{3cm}
    % Asegurar que todos los elementos estén alineados a la izquierda
    \begin{flushleft}
    \languagetag{Python}\\[0.4cm]
    {\fontsize{48}{52}\bfseries\color{primaryColor}Python \color{accentColor}Context\\\color{accentColor}Managers\par}
    \vspace{0.3cm}
    {\fontsize{18}{52}\color{secondaryColor}Elegant Resource Management with the «{\color{accentColor}with}» Statement\par}
    \vspace{0.3cm}
    {\color{secondaryColor}\today\par}
    \vspace{2cm}
    % El QR se alineará a la izquierda aunque internamente esté centrado
    \elegantqr{https://github.com/asanchezyali/social-media-posts/tree/main/Python/ContextManagers}{Source Code}
    \end{flushleft}
}


% En la página del título, aplicar el contenido sin márgenes adicionales
\makeatletter
\renewcommand{\maketitle}{%
    \begin{titlepage}
    \setlength{\parindent}{0pt}
    \setlength{\leftskip}{0pt}
    \titlepagecontents
    \end{titlepage}
}
\makeatother

% También necesitamos corregir la definición para la página final
\newcommand{\finalpagecontents}{%
    \vspace*{4cm}
    \begin{center}
        \begin{tikzpicture}
            % Title text
            \node[text width=14cm, align=center] at (0,0) {
                {\fontsize{48}{52}\bfseries\color{primaryColor}Ready to \color{accentColor}Level Up\\\color{primaryColor}Your Python?\par}
            };
            
            % Add vertical space
            \node at (0,-2) {};
            
            % Question text
            \node[text width=14cm, align=center] at (0,-4) {
                {\fontsize{22}{26}\color{secondaryColor}You've learned something powerful today.\\\vspace{0.5cm}Try creating a context manager in your next project\\\vspace{0.5cm}and watch your code become cleaner and safer!\par}
            };
        \end{tikzpicture}
    \end{center}
}

\begin{document}
\color{primaryColor} % Explicitly set text color to white for the entire document

\begin{titlepage}
    % Eliminamos todos los espacios adicionales para asegurar la alineación
    \titlepagecontents
\end{titlepage}

\section{What Are Context Managers?}

Context managers in Python provide an elegant way to manage resources automatically. Using the \textbf{\textcolor{accentColor}{with}} statement, they handle setup and cleanup operations, ensuring resources are properly acquired and released, even when errors occur.

\begin{itemize}
    \item \textbf{\textcolor{pythonBlue}{Resource Safety:}} Automatic cleanup even during exceptions
    \item \textbf{\textcolor{pythonBlue}{Code Simplicity:}} Eliminates boilerplate try/finally blocks
    \item \textbf{\textcolor{pythonBlue}{Error Handling:}} Built-in mechanism for exception management
\end{itemize}

\subsection{The «with» Statement}

Compare traditional resource management with context managers:

\begin{macterminal}
\begin{lstlisting}
# Without context manager
file = open('data.txt', 'r')
try:
    content = file.read()
finally:
    file.close()

# With context manager
with open('data.txt', 'r') as file:
    content = file.read()
# File is automatically closed
\end{lstlisting}
\end{macterminal}

\section{Creating Context Managers}

\subsection{Class-Based Implementation}

To implement a context manager as a class, you need to define the \textbf{\textcolor{accentColor}{\_\_enter\_\_}} and \textbf{\textcolor{accentColor}{\_\_exit\_\_}} methods:

\begin{itemize}
    \item \textbf{\textcolor{pythonBlue}{\_\_enter\_\_}}: Called when entering the \textbf{with} block - handles setup
    \item \textbf{\textcolor{pythonBlue}{\_\_exit\_\_}}: Called when exiting the \textbf{with} block - handles cleanup
\end{itemize}

\begin{macterminal}
\begin{lstlisting}
class SimpleTimer:
    """A simple context manager that measures execution time."""
    
    def __init__(self, name):
        self.name = name
        self.start_time = None
        
    def __enter__(self):
        """
        Called when entering the -with- block.
        Sets up the timer by recording the start time.
        """
        import time
        print(f"Starting timer: {self.name}")
        self.start_time = time.time()
        return self  # This object will be assigned to the variable after -as-
        
    def __exit__(self, exc_type, exc_value, traceback):
        """
        Called when exiting the -with- block (whether normally or due to exception).
        Calculates and displays the elapsed time.
        
        Parameters:
        - exc_type: Type of exception that occurred (or None if no exception)
        - exc_value: The exception object (or None)
        - traceback: The traceback object (or None)
        """
        import time
        elapsed = time.time() - self.start_time
        
        if exc_type is None:
            # No exception occurred
            print(f"Timer {self.name} finished: {elapsed:.4f} seconds")
        else:
            # An exception occurred
            print(f"Timer {self.name} aborted: {elapsed:.4f} seconds")
            print(f"Exception: {exc_type.__name__}: {exc_value}")
            
        # Return False to let exceptions propagate
        # Return True would suppress the exception
        return False

# Using our custom context manager
with SimpleTimer("calculation") as timer:
    print("Performing calculation...")
    # Simulate some work
    total = sum(i**2 for i in range(1000000))
    print(f"Result: {total}")

# The __exit__ method is called automatically when leaving the with block
print("After the with block")
\end{lstlisting}
\end{macterminal}

\subsection{Function-Based Implementation}

Using the \textbf{\textcolor{accentColor}{@contextmanager}} decorator for a simpler approach:

\begin{macterminal}
\begin{lstlisting}
from contextlib import contextmanager

@contextmanager
def file_manager(filename, mode):
    try:
        # Setup code
        f = open(filename, mode)
        yield f  # Yield the resource
    finally:
        # Cleanup code
        f.close()

# Usage
with file_manager('example.txt', 'r') as f:
    content = f.read()
\end{lstlisting}
\end{macterminal}

\section{Practical Applications}

\subsection{Database Transactions}

A database context manager handles connections and transactions safely. This example explains how it works without requiring database expertise:

\begin{macterminal}
\begin{lstlisting}
import sqlite3

class DatabaseConnection:
    """
    A context manager for database connections.
    - Automatically opens and closes database connections
    - Handles transactions (commit on success, rollback on error)
    """
    def __init__(self, db_name):
        # Store the database name for later use
        self.db_name = db_name  
        self.conn = None
        
    def __enter__(self):
        """
        Called when entering the 'with' block.
        This method:
        1. Opens the database connection
        2. Returns the connection object for use in the 'with' block
        """
        print(f"Opening database connection to {self.db_name}")
        # sqlite3.connect creates a connection to the database
        self.conn = sqlite3.connect(self.db_name)
        return self.conn
        
    def __exit__(self, exc_type, exc_val, exc_tb):
        """
        Called when exiting the 'with' block (in any case).
        This method:
        1. Checks if there was an exception (error)
        2. Rolls back changes if there was an error
        3. Commits changes if everything was successful
        4. Always closes the connection
        
        Parameters:
        - exc_type: The type of exception that occurred (or None if no exception)
        - exc_val: The exception object (or None)
        - exc_tb: The traceback object (or None)
        """
        if self.conn:
            if exc_type:  # If an exception/error occurred in the 'with' block
                print("Error occurred! Rolling back changes...")
                # rollback() undoes all changes made in the current transaction
                self.conn.rollback()
            else:  # If everything went well (no exceptions)
                print("Operations successful! Saving changes...")
                # commit() saves all changes made in the current transaction
                self.conn.commit()
                
            print("Closing database connection")
            # Always close the connection to free up resources
            self.conn.close()
            
        # Return False to let exceptions propagate (be raised)
        # Return True would suppress the exception
        return False

# Simple usage example with an in-memory database (no file needed)
with DatabaseConnection(':memory:') as conn:
    # Create a cursor object to execute SQL commands
    cursor = conn.cursor()
    
    print("Creating a table...")
    # Create a simple table to store people
    cursor.execute('''
    CREATE TABLE people (
        id INTEGER PRIMARY KEY,
        name TEXT,
        age INTEGER
    )''')
    
    print("Adding a person to the database...")
    # Insert a person into the table
    cursor.execute('INSERT INTO people (name, age) VALUES (?, ?)', 
                  ('John', 25))
    
    print("Checking if the person was added...")
    # Verify the person was added
    cursor.execute('SELECT * FROM people')
    result = cursor.fetchone()  # Get the first row
    print(f"Person in database: {result}")
    
# When this block ends:
# 1. The __exit__ method is called automatically
# 2. Changes are committed (saved) if no errors occurred
# 3. The connection is closed
print("Outside the with block - connection is already closed")
\end{lstlisting}
\end{macterminal}

This example shows how context managers eliminate the need to manually open and close database connections, and handle transaction management (commit/rollback) automatically.

\subsection{Temporary State Changes}

Temporarily modify settings or environments:

\begin{macterminal}
\begin{lstlisting}
import os
from contextlib import contextmanager

@contextmanager
def change_directory(path):
    """Temporarily change the working directory."""
    old_dir = os.getcwd()
    try:
        print(f"Changing directory to: {path}")
        os.chdir(path)
        yield  # Give control back to the with block
    finally:
        print(f"Changing back to: {old_dir}")
        os.chdir(old_dir)  # Always restore the original directory

# Code runs with /tmp as working directory
with change_directory("/tmp"):
    print(f"Current directory: {os.getcwd()}")
    # Do some work in the temporary directory
    
# Original directory is restored automatically
print("Back to the original directory")
print(f"Current directory: {os.getcwd()}")
\end{lstlisting}
\end{macterminal}

\section{Advanced Features}

\subsection{Exception Handling}

Context managers can control whether exceptions are propagated or suppressed:

\begin{macterminal}
\begin{lstlisting}
class ExceptionHandler:
    def __enter__(self):
        return self
        
    def __exit__(self, exc_type, exc_val, exc_tb):
        if exc_type is not None:
            print(f"Handled: {exc_type.__name__}: {exc_val}")
            return True  # Suppress the exception
        return False  # Propagate any exception

# Exception will be suppressed
with ExceptionHandler():
    result = 1 / 0
    print("This won't execute")
print("But this will")  # Execution continues
\end{lstlisting}
\end{macterminal}

\subsection{Useful Built-in Context Managers}

Python provides several ready-to-use context managers:

\begin{macterminal}
\begin{lstlisting}
from contextlib import suppress, redirect_stdout
import io

# Suppress specific exceptions
with suppress(FileNotFoundError):
    # No exception if file doesn't exist
    open('non_existent.txt').read()

# Redirect output
f = io.StringIO()
with redirect_stdout(f):
    print("Hello, world!")
output = f.getvalue()  # Contains "Hello, world!"
print(output)
\end{lstlisting}
\end{macterminal}

\section{Performance Measurement}

Create a simple timer context manager:

\begin{macterminal}
\begin{lstlisting}
import time
from contextlib import contextmanager

@contextmanager
def timer(label):
    """
    A context manager that measures and prints the execution time
    of the code inside the 'with' block.
    
    Parameters:
    - label: A descriptive name for the operation being timed
    """
    # 1. Setup phase: Record the start time
    start = time.time()
    try:
        # 2. Yield control to the with-block
        #    Note: We don't yield a value here since we don't need
        #    to expose any object to the with-block
        yield
    finally:
        # 3. Cleanup phase: Calculate and display elapsed time
        #    This runs even if an exception occurs in the with-block
        end = time.time()
        elapsed = end - start
        print(f"{label}: {elapsed:.4f} seconds")

# Example 1: Basic usage
with timer("Processing data"):
    # Time-consuming operation
    time.sleep(0.5)  # Simulate work with a delay
    
# Example 2: Nested timers for profiling different parts of code
with timer("Complete operation"):
    # First subtask
    with timer("Data loading"):
        time.sleep(0.2)  # Simulate loading data
    
    # Second subtask
    with timer("Processing"):
        time.sleep(0.3)  # Simulate processing
    
    # Third subtask
    with timer("Saving results"):
        time.sleep(0.1)  # Simulate saving

# Output will show:
# Processing data: 0.5002 seconds
#  Data loading: 0.2001 seconds
#  Processing: 0.3002 seconds
#  Saving results: 0.1002 seconds
# Complete operation: 0.6009 seconds
\end{lstlisting}
\end{macterminal}

\section{Best Practices}

\begin{itemize}
    \item \textbf{\textcolor{pythonBlue}{Use for Resource Management:}} Files, connections, locks
    \item \textbf{\textcolor{pythonBlue}{Keep Context Managers Focused:}} Single responsibility
    \item \textbf{\textcolor{pythonBlue}{Handle Exceptions Properly:}} Decide when to propagate or suppress
    \item \textbf{\textcolor{pythonBlue}{Document Behavior:}} Especially exception handling
    \item \textbf{\textcolor{pythonBlue}{Test Both Paths:}} Success and exception scenarios
\end{itemize}

\section{Conclusion}

Context managers make Python code cleaner, safer, and more maintainable by automating resource management and error handling. Whether using the built-in managers or creating custom ones, they provide an elegant solution for handling resources properly in all circumstances.

\begin{itemize}
    \item \textbf{\textcolor{pythonBlue}{Automation:}} Eliminates manual resource management
    \item \textbf{\textcolor{pythonBlue}{Safety:}} Ensures resources are properly released
    \item \textbf{\textcolor{pythonBlue}{Flexibility:}} Implement via classes or decorators
\end{itemize}

\section{References}
\begin{itemize}
    \item Python Documentation. \textit{The with statement}. \href{https://docs.python.org/3/reference/compound_stmts.html#the-with-statement}{Link}
    
    \item Real Python. \textit{Context Managers and Python's with Statement}. \href{https://realpython.com/python-with-statement/}{Link}
    
    \item GeeksforGeeks. \textit{Context Manager in Python}. \href{https://www.geeksforgeeks.org/context-manager-in-python/}{Link}
    
    \item DataCamp. \textit{Writing Custom Context Managers in Python}. \href{https://www.datacamp.com/tutorial/writing-custom-context-managers-in-python/}{Link}

    \item Python Documentation. \textit{contextlib — Utilities for with-statement contexts}. \href{https://docs.python.org/3/library/contextlib.html/}{Link}

    \item Python Morsels. \textit{Creating a context manager in Python}. 
    \href{https://www.pythonmorsels.com/creating-a-context-manager//}{Link}

    \item Python Tutorial. \textit{Python Context Manager}. 
    \href{https://www.pythontutorial.net/advanced-python/python-context-managers/}{Link}

    \item Code Rabbit. \textit{Guide To Python Context Managers: Beyond Files}. 
    \href{https://analyticsindiamag.com/deep-tech/guide-to-python-context-managers-beyond-files/}{Link}
    
    \item Every Day Codes \textit{Python Context Managers: Beyond in Depth}. 
    \href{https://everyday.codes/python/python-context-managers-in-depth/}{Link}
    \item Translated, Edited and written in collaboration with AI.
\end{itemize}

\clearpage
\thispagestyle{empty}
\finalpagecontents

\end{document}